% META: {"id":"1SPE-PROBCOND-EX-010","chapitre":"1SPE-PROBA-COND","type_objet":"exercice","capacites":["1SPE-PROBA-COND-C1"],"capacites_codes":["C1"],"methodes":["M1"],"parcours":3,"competences":["calculer","raisonner"],"duree_min":15,"mode_creation":"ex_nihilo","sources_inspiration":[],"parametres_sympy":{},"fichier_tex":"chapitres/1SPE-PROBA-COND/exercices/1SPE-PROBCOND-EX-010.tex","corrige_tex":"chapitres/1SPE-PROBA-COND/corriges/1SPE-PROBCOND-CO-010.tex","status":"approved","origin":"generated","status_history":["generated","audited","accepted","approved"],"release_acceptance":"RELEASE_OWNER_BATCH_ACCEPTANCE","acceptance_closure_digest":"sha256:9b3ccf9a81c5520fbb7e03b7b2d3e7bf2057a02834b6d908bf3be5f49f3b553f","acceptance_receipt_digest":"sha256:4fad86f0282e0fa4e0419cca1ad6379ae7af5a280c04a4a90880f90a1c044242"}
% BEGIN-VERIFY
% from sympy import Rational
% # P(A)=0.6, P(B)=0.5, P(AuB)=0.8
% P_A = Rational(3, 5)
% P_B = Rational(1, 2)
% P_AunionB = Rational(4, 5)
% P_AinterB = P_A + P_B - P_AunionB
% assert P_AinterB == Rational(3, 10)
% P_A_B = P_AinterB / P_A
% assert P_A_B == Rational(1, 2)
% P_B_A = P_AinterB / P_B
% assert P_B_A == Rational(3, 5)
% # P_Abar(B) = P(Abar inter B)/P(Abar)
% P_Abar = 1 - P_A
% P_AbarinterB = P_B - P_AinterB
% assert P_AbarinterB == Rational(1, 5)
% P_Abar_B = P_AbarinterB / P_Abar
% assert P_Abar_B == Rational(1, 2)
% END-VERIFY
\begin{exercice}{1SPE-PROBCOND-EX-010}{3}{15}
On donne $P(A) = \dfrac{3}{5}$, $P(B) = \dfrac{1}{2}$ et $P(A \cup B) = \dfrac{4}{5}$.

\begin{enumerate}
  \item Calculer $P(A \cap B)$.
  \item Calculer $P_A(B)$ et $P_B(A)$.
  \item Calculer $P_{\overline{A}}(B)$.
  \item Comparer $P_A(B)$ et $P_{\overline{A}}(B)$. Interpréter.
\end{enumerate}
\end{exercice}
